% =============================================================================
% Introduction générale
% =============================================================================
\btpLiminary{Introduction générale}

\section*{Contexte}

Le secteur du Bâtiment et des Travaux Publics constitue l'un des piliers de
l'économie camerounaise. Il est caractérisé par une compétition intense où la
réponse aux appels d'offres représente un enjeu stratégique majeur. Chaque
année, des centaines d'entreprises et de bureaux d'études consacrent des
ressources considérables à la préparation de dossiers de candidature, un
processus qui mobilise simultanément des compétences juridiques, techniques,
financières et administratives.

L'émergence des technologies d'intelligence artificielle, et plus
particulièrement des grands modèles de langage, ouvre des perspectives inédites
pour automatiser et optimiser ce processus. Cependant, la spécificité des
appels d'offres — leur caractère hautement structuré, leur exigence de
traçabilité et de conformité, et leur dimension juridique contraignante —
impose une approche réfléchie et méthodique. Un système qui invente une
référence, omet une pièce obligatoire ou cite une clause inexistante ne
constitue pas seulement une gêne~: il expose l'entreprise à l'élimination de
son offre.

\section*{Problématique}

\begin{btpquote}
Comment concevoir et développer une plateforme permettant de capitaliser le
patrimoine documentaire d'un bureau d'études du BTP et de générer
automatiquement des dossiers de réponse aux appels d'offres en s'appuyant sur
ce patrimoine, tout en garantissant la traçabilité des sources, la conformité
réglementaire et la validation humaine des livrables~?
\end{btpquote}

Cette question centrale se décline en plusieurs interrogations
opérationnelles~:

\begin{enumerate}
  \item Comment transformer un patrimoine documentaire hétérogène, non
        structuré et partiellement numérisé en une base de connaissances
        interrogeable~?
  \item Quelle stratégie de recherche d'information permet de retrouver de
        manière fiable à la fois des références terminologiques exactes et des
        concepts métier exprimés différemment~?
  \item Comment contraindre un modèle de langage à produire un texte
        exploitable, sourcé et vérifiable plutôt qu'un texte simplement
        plausible~?
  \item Quelle place accorder à l'opérateur humain pour que le gain de
        productivité ne se paie pas d'une perte de fiabilité~?
\end{enumerate}

\section*{Objectifs}

\subsection*{Objectif général}

Concevoir, développer et valider une plateforme de capitalisation documentaire
et de génération assistée de dossiers de réponse aux appels d'offres, adaptée
au contexte d'un bureau d'études du BTP camerounais.

\subsection*{Objectifs spécifiques}

\begin{itemize}
  \item Caractériser le processus métier de réponse aux appels d'offres et en
        identifier les points de friction documentaires~;
  \item Concevoir une chaîne d'ingestion capable de préserver la structure
        hiérarchique des pièces d'un dossier d'appel d'offres~;
  \item Concevoir un moteur de recherche hybride combinant recherche lexicale
        et recherche sémantique, complété par un ré-ordonnancement fin~;
  \item Concevoir un mécanisme de génération contrainte produisant des sections
        structurées, citées et traçables~;
  \item Intégrer un dispositif de supervision humaine à chaque étape critique~;
  \item Évaluer empiriquement la performance du système sur des indicateurs de
        recherche et de génération.
\end{itemize}

\section*{Hypothèses de recherche}

\begin{btptable}{Hypothèses de recherche et critères de validation}
\begin{tabularx}{\linewidth}{@{}C{1.1cm} Y L{3.5cm}@{}}
\toprule
\rowcolor{btppale}
\btpth{Réf.} & \btpth{Hypothèse} & \btpth{Critère de validation} \\
\midrule
H1 & La recherche hybride surpasse la recherche sémantique seule sur un corpus d'appels d'offres. & Gain de rappel à 20 documents \\
H2 & Le découpage hiérarchique parent-enfant améliore la qualité des réponses par rapport à un découpage à fenêtre fixe. & Gain d'ancrage et d'exhaustivité \\
H3 & Le ré-ordonnancement par modèle croisé surpasse la seule fusion de rangs. & Gain de précision à 5 documents \\
H4 & La génération section par section améliore la conformité par rapport à une génération globale. & Gain de couverture des exigences \\
H5 & La supervision humaine améliore la fiabilité des livrables sans annuler le gain de productivité. & Taux de validation et temps de cycle \\
\bottomrule
\end{tabularx}
\end{btptable}

\section*{Organisation du mémoire}

Ce mémoire s'organise en trois chapitres. Le \textbf{chapitre~1} dresse l'état
de l'art~: il caractérise le domaine des appels d'offres dans le BTP, puis
parcourt les disciplines mobilisées, de la gestion documentaire à la génération
augmentée par récupération, en passant par la recherche d'information, la
compréhension de documents structurés et la résolution d'entités. Le
\textbf{chapitre~2} expose la méthodologie~: analyse des besoins, modélisation
des données, architecture du système, stratégies de recherche et de génération,
protocole d'évaluation et planification. Le \textbf{chapitre~3} présente la
plateforme développée, son architecture technique, les résultats de
l'évaluation et une analyse comparative des approches candidates. Une
conclusion générale récapitule les contributions et ouvre sur les perspectives.

\btpEndMark
\clearpage
